\scanpage{130}
\begin{center}- \mspage{117} -\end{center}\margin{IX}

Pour tout entier $n$, l'homomorphisme correspondant $X(n)\longrightarrow Q(n)$, étant surjectif sur les fibres génériques, est surjectif ; comme $Q(n)$ est étale sur $\unsure{n|\eta}$, donc $X(n)\longrightarrow Q(n)$ est plat, il s'ensuit que c'est un épimorphisme. Donc $X\longrightarrow Q$ \uline{est un épimorphisme composant par composant}. Si donc $\bar F$ désigne le noyau de $X\longrightarrow Q$, il s'ensuit que $\bar F$ est un groupe de Barsotti-Tate et $X$ est une extension de $Q$ par $\bar F$. D'ailleurs l'isomorphisme $F_\eta\overset\sim\longrightarrow\bar F_\eta$ se prolonge, par le théorème de TATE, en un isomorphisme $F\simeq\bar F$. Par suite, $X$ apparaît comme une extension de $Q$ par $F$. Considérons alors le quotient $\bar E=X/P$, on voit encore, par application du théorème de TATE, qu'on a un isomorphisme $E\simeq\bar E$. Il s'ensuit que $X$ est une extension panachée de $E$ par $F$ prolongeant l'extension panachée $X_\eta$, cqfd.

\uline{Remarques} 9.5.6.4. a) Bien entendu, la dernière assertion de 9.5.6 devrait être valable sans restriction sur $K$, puisqu'il devrait en être ainsi du théorème de Tate invoqué.

b) On peut donner une autre démonstration du critère de bonne réduction 5.10, en se ramenant comme dans 5.10 au cas où $A_K$ est à réduction semi-stable, et en appliquant alors 9.5.6 à $T_\ell(A_K)$ considérée comme extension panachée (cf. 9.6), qui implique $u_\ell=0$, et utilisant 11.6.2 a) plus bas qui implique que l'on a $T_\ell(A_K)^t=0$.

9.5.8. Lorsque $\ell\neq p$, les groupes de Barsotti-Tate envisagés sur $S$ resp. sur $K$ s'interprètent simplement en termes de faisceaux $\ell$-adiques constants tordus sur $S$, ou de façon équivalente, en termes de représentations de $\pi_{\unsure{0|o}}=\operatorname{Gal}(\bar k/k)$ resp. de $\pi=\operatorname{Gal}(\bar K/K)$ sur des $\mathbb Z_\ell$-modules libres

\margin{435}\margin{129}
